\documentclass[10pt,a4paper]{amsart}
\usepackage[margin=19mm]{geometry}
\usepackage{amsmath,amssymb,mathtools}
\usepackage[scr=rsfs,cal=euler]{mathalfa}
\usepackage{xfrac}
\usepackage[unicode,hidelinks,bookmarks=false]{hyperref}
\newcommand{\ie}{i.e.\ }
\newcommand{\cf}{cf.\ }
\newcommand{\resp}{respectively\ }
\newcommand{\Subsection}{\subsection}
\providecommand{\nobreakdash}{\nobreak\hbox{-}}
\setlength{\emergencystretch}{2em}
\multlinegap0pt
\usepackage{amssymb}
\usepackage{enumerate}
\usepackage{dsfont}
\usepackage{mathtools}
\usepackage{cancel}
\usepackage{float}
\newtheorem{lemma}{Lemma}
\newtheorem{sublemma}{Sublemma}
\newcommand{\QA}[1]{(\forall{#1})\,}
\newcommand{\QE}[1]{(\exists{#1})\,}
\title{Intuitionism and computing with partial information}
\author{Hristo Ganchev, Paul Shafer, Theodore A. Slaman, Andrea Sorbi, Mariya I. Soskova}
\date{Source: arXiv:2605.20841v1 (CC BY 4.0)}
\begin{document}
\maketitle

\noindent\emph{Source and licence.} Intuitionism and computing with partial information. Version arXiv:2605.20841v1 (CC BY 4.0), distributed by arXiv under the Creative Commons Attribution 4.0 International licence, http://creativecommons.org/licenses/by/4.0/, as recorded in the arXiv licence metadata for this exact version and shown on the abstract page https://arxiv.org/abs/2605.20841v1. Full licence notice: \texttt{licenses/arxiv-2605.20841-CC-BY-4.0.txt}.\par\smallskip

\noindent\textit{Editorial scope.} Counted unit: the article's lemma lem:a-splitting (source0.tex lines 1658--1679). The article's complete original proof of it is delivered with it (source0.tex lines 1681--1868, one \texttt{proof} environment of 188 lines). No mathematics is written, repaired or re-derived: the statement and the proof are the article's own lines, in the article's own notation, and no retained line is substituted or re-flowed.  No further source-local context is needed: every cross-reference of every retained line resolves inside the two delivered spans.  Nothing was omitted from any retained span, so the package carries no omission record.  No retained line contains a citation, so the wrapper carries no bibliography and no bibliography entry.  No retained line contains a figure environment, an image-inclusion command or drawing code, so no image is delivered and the package declares no asset dependency.  Editorial formatting only, in addition: an AMS \texttt{amsart} preamble that restates the article's own declarations and macros used by the retained lines, namely the environments \texttt{lemma}, \texttt{sublemma} on separate counters, together with the standard packages those lines need.  The retained environments print in this document as Lemma 1, Sublemma 1 in order of appearance, whereas the article prints the same items with the numbering of its own full document.  The complete native source of the article as distributed in the arXiv e-print for this exact version is bundled unchanged as \texttt{sources/201022/source0.tex}.
\begin{lemma}\label{lem:a-splitting}
Suppose that $\mathcal{I},\mathcal{G}$ are sets of sets of numbers and 
that $A, B$ are sets such that 

\begin{enumerate}
  \item[(a)] $\mathcal{I}$ is countable and downwards~$\leq_e$-closed;
  \item[(b)] $\mathcal{G} \subseteq \mathcal{I}^c$ is countable and closed 
      under $\equiv_e$; 
  \item[(c)] $A, B \in \mathcal{I}$ and $B \nleq_e A$. 
\end{enumerate} 

Then there exists a generalized enumeration operator $\Gamma$ satisfying 

\begin{enumerate}
  \item[(0)] $\Gamma \oplus A \nleq_e A$
  \item if $X \in \mathcal{I}$ then $\Gamma \oplus A\geq_e X \textrm{ if 
      and only if } A \geq_e X$; 
  \item if $X \in \mathcal{G}$ then $\Gamma \oplus A \ngeq_e X$; 
  \item $\Gamma(B) \notin \mathcal{I}$;
  \item $\Gamma\oplus A \ngeq_e \Gamma(B)$.
\end{enumerate}
\end{lemma}

\begin{proof}
Obtain $\Gamma$ by forcing with $\langle P, \leq_P \rangle$ where 
\begin{enumerate}
  \item[I.] each condition $p \in P$ is a triple $p=\langle \Gamma_p, B_p, 
      N_p\rangle$, where 
      
\begin{enumerate}
  \item[-] $\Gamma_p$ is a finite set of axioms $\langle n, u\rangle$ 
      where $n\in \omega$ and $u$ is the canonical index of the finite 
      set $D_u$; 
  \item[-] $B_p$ is a finite subset of $B$;
  \item[-] $N_p$ is a finite subset of $\omega$.
\end{enumerate}
  \item[II.] $p\leq_P q$ if
  \begin{enumerate}
    \item[-] $\Gamma_q \subseteq \Gamma_p$ (each axiom of $\Gamma_q$ is 
        an axiom of $\Gamma_p$);
    \item[-] $B_q \subseteq B_p$ and $N_q \subseteq N_p$;
    \item[-] $\QA{\langle n, u\rangle \in \Gamma_p \smallsetminus 
        \Gamma_q} \left[ B_p \subseteq D_u \right]$ (the new axioms must 
        use $B_p$); 
    \item[-] $\QA{\langle n, u\rangle \in \Gamma_p \smallsetminus 
        \Gamma_q} \left[ n \notin N_p \right]$ (no new axiom enumerates 
        elements of $N_p$). 
  \end{enumerate}
\end{enumerate} 

We will exhibit appropriate dense sets, so that if $\Gamma$ comes from a 
sufficiently generic filter then it satisfies the claims of the lemma.

The following sublemma addresses item (0) of Lemma~\ref{lem:a-splitting}. 

\begin{sublemma}\label{lem:0}
$\Gamma \nleq_e A$.
\end{sublemma}

\begin{proof}[Proof~of~Sublemma~\ref{lem:0}]
Let $\Theta$ be an $e$-operator. The set of conditions 
\[ 
\mathcal{C}=\{p: \QE{m}\left[m \in \Theta(A) \,\&\, p \Vdash m 
\notin \Gamma\right]\} \cup
\{\QE{m}\left[m \notin \Theta(A) \,\&\, p \Vdash m \in \Gamma\right]\}
\]
is dense. To see this, suppose $p=\langle \Gamma_p, B_p, N_p\rangle$ is a 
condition. Let $m=\langle m_0, m_1\rangle$ be such that $m, m_0, m_1$ are 
larger than any number mentioned in $p$, with $D_{m_1} \supseteq B_p$.
If $m \in \Theta(A)$, then let
$q=\langle \Gamma_p, B_p, N_p\cup \{m_0\}\rangle$
so that $q\Vdash \langle m_0, m_1\rangle \notin \Gamma$.
If $m \notin \Theta(A)$, then let $q=\langle \Gamma_p \cup \{m\}, B_p, N_p\rangle$.
In either case there is a condition $q \leq_P p$ such that $q \in \mathcal{C}$. 
\end{proof}

The following sublemma addresses items (1) and (2) of 
Lemma~\ref{lem:a-splitting}. 
 
\begin{sublemma}\label{lem:1}
For $X \in \mathcal{I}$,  $\Gamma \oplus A\geq_e X $ if and only if $A \geq_e 
X$. For $X \in \mathcal{G}$, $\Gamma \oplus A \ngeq_e X$. 
\end{sublemma}

\begin{proof}[Proof~of~Sublemma~\ref{lem:1}] 
Let $X \in \mathcal{I} \cup \mathcal{G}$. Suppose that $\Theta$ is an $e$-operator and 
that $\Theta(\Gamma\oplus A)=X$. This implies that the set of conditions 
\begin{align*}
\mathcal{C} &=\{p: p \Vdash \Theta(\Gamma\oplus A)\ne X\}\\
            &=\{p: \QE{m}\left [m \notin X \,\&\, p \Vdash m \in \Theta(\Gamma 
      \oplus A)\right]\} \cup \{p: \QE{m} \left[ m \in X \,\&\, p \Vdash m \notin 
      \Theta(\Gamma \oplus A) \right]\} 
\end{align*}
is not dense. (Here, if $p=\langle \Gamma_p, B_p, N_p \rangle$ is a 
condition, then $p \Vdash m \in \Theta(\Gamma \oplus A)$ means $m \in 
\Theta(\Gamma_p \oplus A)$, and $p \Vdash m \notin \Theta(\Gamma \oplus A)$ 
means that $\QA{q \le_{P} p}[m \notin \Theta(\Gamma_q \oplus A)]$. Analogous 
definitions will be tacitly assumed throughout the proof of 
Lemma~\ref{lem:a-splitting}.) Hence there exists a condition $p$ with no 
extension in this set, which implies 
\[
\QA{m} \big[ m \in X \Leftrightarrow \QE{q\leq_P p} 
[m \in \Theta(\Gamma_q \oplus A) ] \big].
\] 
Then $A\geq_e X$: indeed, $A$ enumerates $m \in X$ by searching for a 
condition $q \leq_P p$ forcing $m \in \Theta(\Gamma_q\oplus A)$. The set of 
such conditions $q$ is positively c.e. in $A$. So we conclude that 
\[
\QA{X \in \mathcal{I} \cup \mathcal{G}}[ \Gamma \oplus A \geq_e
X \Leftrightarrow A \geq_e X]. \qedhere
\] 
\end{proof}

The next sublemma addresses item (3) of Lemma~\ref{lem:a-splitting}. 

\begin{sublemma}\label{lem:2}
$\Gamma(B) \notin \mathcal{I}$. 
\end{sublemma}

\begin{proof}[Proof~of~Sublemma~\ref{lem:2}] 
For every condition $p=\langle \Gamma_p, B_p, N_p\rangle$ and $X \in 
\mathcal{I}$, condition $p$ has an extension in the set $\mathcal{C}$ of 
conditions 
\[
\mathcal{C}=\{q: \QE{m} \left[m \in X \,\&\, 
q \Vdash m \notin \Gamma(B)\right]\}
\cup \{q: \QE{m} \left[m \notin X \,\&\, q \Vdash m \in \Gamma(B)\right]\},
\]
which therefore is dense, and thus $X \ne \Gamma(B)$. Indeed, to see density, 
let $p=\langle \Gamma_p, B_p, N_p\rangle$ be a condition, and let $m$ be 
larger than any number mentioned in $p$. If $m\in X$, then the condition 
$q=\langle \Gamma_p, B_p, N_p \cup \{m\} \rangle$ extends $p$ and forces $m 
\notin \Gamma(B)$. If $m\notin X$, then the condition $q=\langle \Gamma_p\cup 
\{\langle m, u\rangle\}, B_p, N_p  \rangle$ where $D_u = B_p$ extends $p$ with $q \Vdash m \in 
\Gamma(B)$. 
\end{proof}

The next sublemma addresses item (4) of Lemma~\ref{lem:a-splitting}. 

\begin{sublemma}\label{lem:3}
$\Gamma\oplus A \ngeq_e \Gamma(B)$.
\end{sublemma}

\begin{proof}[Proof~of~Sublemma~\ref{lem:3}] 
Suppose that $p=\langle \Gamma_p, B_p, N_p\rangle$ is a condition, and let 
$\Theta$ be an $e$-operator. Let $\mathcal{C}$ be the set of conditions 
\begin{multline*}
\mathcal{C}=\{q: \QE{m} \left[ q \Vdash m \in \Theta(\Gamma \oplus A) \,\&\,
q \Vdash m \notin \Gamma(B)\right]\},\\
\cup \{q: \QE{m} \left[q \Vdash m \in \Gamma(B) \,\&\, q \Vdash m \notin 
\Theta(\Gamma \oplus A) \right]\}.
\end{multline*}
In Case~1 and Case~2 below, $m$ and $d$ vary through the numbers greater than 
any number mentioned in $p$. 

\underline{Case~1.} $\QE{m}\QE{d \notin B}\QE{q \leq_P p}
\bigl[m \in \Theta(\Gamma_q\oplus A) \,\&\, \QA{\langle n, u 
\rangle \in \Gamma_q \smallsetminus \Gamma_p}[d \in D_u]\bigr]$.

Let
\[
r=\langle \Gamma_q, B_q,  N_q \cup\{m\}\rangle.
\]
Then $r <_P p$, $r\Vdash   m \notin \Gamma(B)$ since $\Gamma_p$ does not 
mention $m$ and the axioms in $\Gamma_q \smallsetminus \Gamma_p$ do not apply 
to $B$, and $ r \Vdash m \in \Theta(\Gamma \oplus A)$. Hence $p$ has an 
extension $r \in \mathcal{C}$. 

\underline{Case~2.} $\QE{m}\QE{d \in B}[ \langle \Gamma_p, B_p\cup \{d\}, 
N_p\rangle \Vdash m \notin \Theta(\Gamma \oplus A)]$ (thus requiring that 
using $d$ in all new axioms makes it impossible for $\Theta(\Gamma \oplus A)$ 
to enumerate $m$). 

Let $r=\langle \Gamma_p \cup \{\langle m, u \rangle\}, 
B_p\cup\{d\}, N_p \rangle$, where $D_u = B_p\cup \{d\}$. Then $r<_P p$, 
$r\Vdash m \notin \Theta(\Gamma 
\oplus A)$, and $r \Vdash m \in \Gamma(B)$. 

Again, $p$ has an extension $r \in \mathcal{C}$.

\underline{Case~3.} Otherwise. 

Then, for all sufficiently large $d$, 
\begin{equation}\label{eq:case3}
d \in B \Leftrightarrow\\
\QE{m}\QE{q \leq_P p}
\big[m \in \Theta (\Gamma_q \oplus A)\;\&\; \QA{\langle n, u\rangle \in 
\Gamma_q \smallsetminus \Gamma_p} 
    [d \in D_u]\big].
\end{equation}
Indeed, the righthand side of (\ref{eq:case3}) implies that $d\in B$, since 
otherwise Case~1 would hold. On the other hand, if $d \in B$, then 
$\langle\Gamma_p, B_p\cup \{d\}, N_p\rangle \nVdash m \notin \Theta(\Gamma 
\oplus A)$ for every sufficiently large $m$ since Case~2 does not hold. This 
means that there is a $q \leq_P \langle\Gamma_p, B_p\cup \{d\}, N_p\rangle 
\leq_P p$ with $m \in \Theta(\Gamma_q \oplus A)$, and $\QA{\langle n, 
u\rangle \in \Gamma_q \smallsetminus \Gamma_p}[d \in D_u]$ by the definition 
of extension.  Thus the righthand side of (\ref{eq:case3}) holds. 

But then, contradicting the assumption, it follows that $B\leq_e A$ since the 
description in the righthand side of (\ref{eq:case3}) is positive in $A$. 

It follows that every $p$ has an extension in $\mathcal{C}$, so $\mathcal{C}$ 
is a dense set of conditions. 
\end{proof}

In conclusion, if $\Gamma$ is sufficiently generic and intersects the dense 
sets of conditions exhibited in the proof then the generalized enumeration 
operator $\Gamma$ has the desired properties, which satisfy 
Lemma~\ref{lem:a-splitting}. 
\end{proof}

\end{document}
